\documentclass[a4paper,10pt]{article}
\usepackage[utf8]{inputenc} 
\usepackage[english]{babel}
\usepackage[scale=.7]{geometry}
\usepackage{graphicx}
\usepackage{amsmath,amssymb,amsthm}
\usepackage{hyperref}
\usepackage{polynomial}
\usepackage{enumitem}
\usepackage{hyperref}

\numberwithin{equation}{section} 
\newtheorem{theorem}[subsubsection]{Theorem}
\newtheorem{corollary}[subsection]{Corollary}
\newtheorem{lemma}[subsection]{Lemma}
\newtheorem{proposition}[subsection]{Proposition}
\theoremstyle{definition}
\newtheorem{definition}[subsubsection]{Definition}
\newtheorem{remark}[subsection]{Remark}
\newtheorem{example}[subsection]{Example}

\title{The Cayley-Hamilton Theorem: A Pillar of Linear Algebra}

\author{Johan Voskamp}

\date{\today}

\begin{document}
\maketitle

\begin{abstract}
    temp
\end{abstract}
\section{AI/Data usage}
\section{Introduction}
\section{Prerequisites}
    To understand parts of the disproof of the - Theorem, we will have to define and explain some fundamental concepts beforehand.
    
    \subsection{Cayley Graphs}
    \subsection{Minimal Cayley Graphs}
       Note that for the generating set S, if it is minimal
    \subsection{No-lonely-colored graphs}
        
    \subsection{Temp}

\section{a}
   

\section{b}

\bibliography{Bibliography}
\bibliographystyle{plain}

\end{document}
